\chapter*{About this guide}
\addcontentsline{toc}{chapter}{About this guide}
\markboth{About this guide}{}

Calango is a desktop application for building, visualizing, simulating and
analysing atomic structures. It pairs a C++20 / OpenGL core --- fast enough
to orbit tens of thousands of atoms in real time --- with an embedded
CPython interpreter running the Atomic Simulation Environment (ASE), so
every structure you see on screen can be handed to a real calculator
without leaving the program.

This guide is written for the person using Calango, not the person
building it. It walks through the workspace, then covers each menu in the
order you are likely to need it: getting data in, building and editing
structures, making them look right, running simulations locally or on a
cluster, and analysing the results. Chapter~\ref{ch:reference} is a
reference you can jump straight to --- keyboard shortcuts, the complete
menu map, file formats and troubleshooting.

\section*{Conventions}

\begin{description}
  \item[\menu{Analysis}\menusep\menu{Raman Modes}] A menu path. Each
    \menusep{} step is one level deeper.
  \item[\ui{Compute}] A button, field, checkbox or tab label exactly as it
    appears in the interface.
  \item[\kbd{Ctrl}+\kbd{R}] A keyboard shortcut. On macOS, \kbd{Ctrl}
    means \kbd{Cmd} wherever the shortcut is a standard editing action
    (open, save, quit, undo); single-letter viewport shortcuts have no
    modifier on any platform.
  \item[\file{run.py}] A file, directory or path.
\end{description}

Three kinds of callout appear throughout:

\begin{calnote}
Background you do not strictly need, but that explains why something
behaves the way it does.
\end{calnote}

\begin{caltip}
A shortcut or working habit that saves time.
\end{caltip}

\begin{calwarn}
An external dependency --- a Python package, a solver binary, a network
connection --- that the feature cannot work without.
\end{calwarn}

\Needspace{6\baselineskip}
\section*{Version}

This guide documents Calango \textbf{26.7.83}. The version you are running
is shown in \menu{Help}\menusep\menu{About Calango}, together with the
Python and ASE versions actually in use --- worth checking first whenever
something behaves unexpectedly.
